\section{Metodología experimental}
\subsection{Diseño factorial y control de entradas}
Se evaluaron tres tamaños, dos políticas, tres perfiles de demanda y dos condiciones de fallo. Cada combinación tuvo 30 réplicas: $3\times2\times3\times2\times30=1.080$ ejecuciones.

\begin{table}[t]
\caption{Parámetros del experimento principal}
\label{tab:parametros}
\centering\footnotesize
\begin{tabularx}{\columnwidth}{@{}p{.36\columnwidth}Y@{}}\toprule
Parámetro & Valor\\\midrule
Infraestructura & 4 hosts, 8 núcleos y 16 GiB por host\\
VMs & 4, 6 u 8; 2 vCPU y 4 GiB por VM\\
Cola por VM & FIFO; 100 solicitudes esperando\\
Servicio por puesto & Exponencial; media 0,10 s\\
Retardo de red & 0,02 s fijo, sin ocupar CPU\\
Objetivo / vencimiento & 2 s / 10 s desde la llegada\\
Calentamiento & 600 s con la tasa inicial\\
Medición / seguimiento & 3.600 s / 10 s\\
Fallo / detección & Host 1: 120 s / 5 s\\
Réplicas & 30 por combinación\\\bottomrule
\end{tabularx}
\end{table}

La carga baja mantuvo 40 solicitudes/s y la alta, 72. El perfil con pico utilizó 40, 120 y 40 solicitudes/s en tres tramos de 1.200 s. Cada perfil se ejecutó sin fallo y con caída del host 1 entre los segundos 1.800 y 1.920 de medición. Las capacidades nominales sin fallo fueron 80, 120 y 160 solicitudes/s; con fallo, 60, 80 y 120. Una intensidad nominal igual a uno no asegura buen servicio durante el intervalo afectado.

Se utilizaron dos flujos aleatorios separados para llegadas y servicio, derivados de \texttt{SeedSequence} de NumPy con base 20261029, índice de réplica e índice de carga. Cada solicitud recibió su trabajo antes de aplicar el balanceo. Para una misma carga y réplica se reutilizó la traza completa en todas las alternativas y condiciones de fallo; su identidad se registró mediante SHA-256. Este acoplamiento evita atribuir a la política diferencias causadas por entradas distintas.

La medición por solicitudes incluye las llegadas durante la ventana observada, con seguimiento hasta su terminación o vencimiento. Las solicitudes del calentamiento no forman parte de esa cohorte, aunque sí pueden ocupar recursos al iniciar la ventana. Para el costo se usó un horizonte común de $T=(600+3.600+10)/3.600=1,16944$ horas, sin asumir que una hora simulada equivale a una hora de ejecución del programa.

\subsection{Indicadores e incertidumbre}
Se registraron $A$, $G$, latencia media y p95 de respuestas correctas, espera media de esas respuestas, longitud de cola ponderada por tiempo y utilización. La cola media del sistema es:
\begin{equation}
\overline Q=\frac{1}{H}\int_0^H\sum_{i=1}^{c}Q_i(t)\,dt,
\end{equation}
donde $H=3.600$ s. La utilización es el tiempo de vCPU ocupado dividido entre $2cH$; incluye la capacidad contratada durante los fallos en el denominador. Las integrales se actualizan exactamente en cada cambio de estado, no mediante el muestreo de las figuras.

El throughput contabiliza respuestas correctas finalizadas dentro de la ventana dividido entre $H$; difiere de la proporción de éxito de la cohorte, que contempla el seguimiento. El p95 se calcula por réplica sobre respuestas correctas; los errores se reflejan en $A$ y $G$. Por ello no se utiliza una baja latencia condicionada a éxitos como evidencia suficiente de buen servicio.

Para cada indicador se calculó un intervalo puntual del 95\,\%:
\begin{equation}
\bar y\pm t_{0,975;29}\frac{s_y}{\sqrt{30}}.
\end{equation}
El p95 informado es el promedio de los p95 de las réplicas. Las diferencias MP--RR se analizaron por semilla mediante intervalos pareados. No se trataron las solicitudes de una misma ejecución como observaciones independientes para estimar la incertidumbre entre réplicas. Los intervalos son aproximados y no están ajustados para comparaciones múltiples.

\subsection{Verificación y contraste analítico}
Se distinguió la verificación del programa del alcance de su validación, siguiendo la separación conceptual discutida por Sargent \cite{sargent}. Se implementaron diez pruebas deterministas para comprobar FIFO, capacidad de cola, liberación de CPU por vencimiento, eliminación de esperas vencidas, fallo y detección, retardo sin retener CPU, cohorte y seguimiento, integrales de estado, selección MP y reproducibilidad de trazas.

Cada ejecución comprobó además la identidad entre solicitudes ofrecidas y sus cuatro estados terminales excluyentes: correctas, rechazadas, vencidas y fallidas. Se verificó que ninguna VM excediera dos solicitudes simultáneas y que no quedaran solicitudes activas al terminar el seguimiento.

El contraste analítico utilizó M/M/1 con $\lambda=5$, $\mu=10$, un puesto, cola ilimitada, sin red, plazos ni fallos. Se realizaron 30 réplicas adicionales con 1.000 s de calentamiento y 10.000 s de medición. Las referencias son $W=0,20$ s, $L=1$ y $\rho=0,5$. Esta prueba acredita coherencia con un caso conocido; no valida el modelo completo contra un centro de datos real.
